% from ega1-10-fr.tex
\oldpage[I]{195}

$\varprojlim_{\Phi}
(\mathscr{F}\otimes_{\mathscr{O}_X}
(\mathscr{O}_X/\mathscr{J}))$~; on dit que ses sections au-dessus de $X'$ sont
les \emph{sections formelles de $\mathscr{F}$ le long de $X'$}.
\end{definition}

Il est immédiat que pour tout ouvert $U\subset X$, on a
$(\mathscr{F}|U)_{/(U\cap X')}=(\mathscr{F}_{/X'})|(U\cap X')$.

Par passage à la limite projective, il est clair que
$(\mathscr{O}_X)_{/X'}$ est un faisceau d'anneaux, et que
$\mathscr{F}_{/X'}$ peut être considéré comme un
$(\mathscr{O}_X)_{/X'}$-Module. En outre, comme il existe une base de la
topologie de $X'$ formée d'ouverts quasi-compacts, on peut considérer
$(\mathscr{O}_X)_{/X'}$ (resp. $\mathscr{F}_{/X'}$) comme un
\emph{faisceau d'anneaux topologiques} (resp. de
\emph{groupes topologiques}) limite projective des faisceaux d'anneaux
(resp. groupes) \emph{pseudo-discrets}
$\mathscr{O}_X/\mathscr{J}$ (resp.
$\mathscr{F}\otimes_{\mathscr{O}_X}
(\mathscr{O}_X/\mathscr{J})=\mathscr{F}/\mathscr{J}\mathscr{F}$), et, par
passage à la limite projective, $\mathscr{F}_{/X'}$ devient alors un
\emph{$(\mathscr{O}_X)_{/X'}$-Module topologique}
(\hyperref[0.3.8.1-fr]{0, 3.8.1} et
\hyperref[0.3.8.2-fr]{3.8.2})~; rappelons que pour tout ouvert
\emph{quasi-compact} $U\subset X$,
$\Gamma(U\cap X',(\mathscr{O}_X)_{/X'})$ (resp.
$\Gamma(U\cap X',\mathscr{F}_{/X'})$) est alors limite projective des
anneaux (resp. groupes) discrets $\Gamma(U,\mathscr{O}_X/\mathscr{J})$
(resp. $\Gamma(U,\mathscr{F}/\mathscr{J}\mathscr{F})$).

Si maintenant $u:\mathscr{F}\to\mathscr{G}$ est un homomorphisme de
$\mathscr{O}_X$-Modules, on en déduit canoniquement des homomorphismes
$u_{\mathscr{J}}:
\mathscr{F}\otimes_{\mathscr{O}_X}(\mathscr{O}_X/\mathscr{J})
\to
\mathscr{G}\otimes_{\mathscr{O}_X}(\mathscr{O}_X/\mathscr{J})$
pour tout $\mathscr{J}\in\Phi$, et ces homomorphismes forment un système
projectif. Par passage à la limite projective et restriction à $X'$, ils
donnent donc un $(\mathscr{O}_X)_{/X'}$-homomorphisme continu
$\mathscr{F}_{/X'}\to\mathscr{G}_{/X'}$, noté $u_{/X'}$ ou $\widehat u$,
et appelé le \emph{complété de l'homomorphisme $u$ le long de $X'$}. Il est
clair que si $v:\mathscr{G}\to\mathscr{H}$ est un second homomorphisme de
$\mathscr{O}_X$-Modules, on a
$(v\circ u)_{/X'}=(v_{/X'})\circ(u_{/X'})$ donc
$\mathscr{F}_{/X'}$ est un \emph{foncteur additif covariant} en
$\mathscr{F}$, de la catégorie des $\mathscr{O}_X$-Modules cohérents, à
valeurs dans la catégorie des $(\mathscr{O}_X)_{/X'}$-Modules topologiques.

\begin{proposition}[10.8.5]
\label{I.10.8.5-fr}
Le support de $(\mathscr{O}_X)_{/X'}$ est $X'$~; l'espace topologiquement
annelé $(X',(\mathscr{O}_X)_{/X'})$ est un préschéma formel localement
noethérien, et si $\mathscr{J}\in\Phi$, $\mathscr{J}_{/X'}$ est un Idéal de
définition de ce préschéma formel. Si $X=\operatorname{Spec}(A)$ est un
schéma affine d'anneau noethérien, $\mathscr{J}=\widetilde{\mathfrak{J}}$
où $\mathfrak{J}$ est un idéal de $A$, et $X'=V(\mathfrak{J})$,
$(X',(\mathscr{O}_X)_{/X'})$ s'identifie canoniquement à
$\operatorname{Spf}(\widehat A)$, où $\widehat A$ est le séparé complété de
$A$ pour la topologie $\mathfrak{J}$-préadique.
\end{proposition}

\begin{proof}
On peut évidemment se borner à prouver la dernière assertion. On sait
(\hyperref[0.7.3.3-fr]{0, 7.3.3}) que le séparé complété
$\widehat{\mathfrak{J}}$ de $\mathfrak{J}$ pour la topologie
$\mathfrak{J}$-préadique s'identifie à l'idéal
$\mathfrak{J}\widehat A$ de $\widehat A$, et que $\widehat A$ est un anneau
$\widehat{\mathfrak{J}}$-adique noethérien tel que
$\widehat A/\widehat{\mathfrak{J}}^n=A/\mathfrak{J}^n$
(\hyperref[0.7.2.6-fr]{0, 7.2.6}). Cette dernière relation montre que les
idéaux premiers ouverts de $\widehat A$ sont les idéaux
$\widehat{\mathfrak p}=\mathfrak p\widehat A$, où $\mathfrak p$ est un
idéal premier de $A$ contenant $\mathfrak{J}$, et que l'on a
$\widehat{\mathfrak p}\cap A=\mathfrak p$, d'où
$\operatorname{Spf}(\widehat A)=X'$. Comme
$\mathscr{O}_X/\mathscr{J}^n=(A/\mathfrak{J}^n)^{\sim}$, la proposition
découle aussitôt des définitions.
\end{proof}

On dit que le préschéma formel ainsi défini est le
\emph{complété de $X$ le long de $X'$} et on le note $X_{/X'}$ ou
$\widehat X$ si aucune confusion n'est à craindre. Lorsque l'on prend
$X'=X$, on peut prendre $\mathscr{J}=0$, et on a donc $X_{/X}=X$.

Il est clair que si $U$ est un sous-préschéma induit sur un ouvert de $X$,
$U_{/(U\cap X')}$ s'identifie canoniquement au sous-préschéma formel induit
par $X_{/X'}$ sur l'ouvert $U\cap X'$ de $X'$.

\begin{corollary}[10.8.6]
\label{I.10.8.6-fr}
Le préschéma (usuel) $\widehat X_{\mathrm{red}}$ est l'unique
sous-préschéma réduit de $X$ ayant pour espace sous-jacent $X'$
(\hyperref[I.5.2.1-fr]{5.2.1}). Pour que $\widehat X$ soit noethérien, il
faut et il suffit que $\widehat X_{\mathrm{red}}$ le soit, et il suffit que
$X$ le soit.
\end{corollary}

